\documentclass[10pt,a4paper]{amsart}
\usepackage[margin=19mm]{geometry}
\usepackage{amsmath,amssymb,mathtools}
\usepackage[scr=rsfs,cal=euler]{mathalfa}
\usepackage{xfrac}
\usepackage[unicode,hidelinks,bookmarks=false]{hyperref}
\newcommand{\ie}{i.e.\ }
\newcommand{\cf}{cf.\ }
\newcommand{\resp}{respectively\ }
\newcommand{\Subsection}{\subsection}
\providecommand{\nobreakdash}{\nobreak\hbox{-}}
\setlength{\emergencystretch}{2em}
\multlinegap0pt
\usepackage{mathtools}
\usepackage{amssymb}
\newtheorem{lemma}{Lemma}
\newtheorem{theorem}{Theorem}
\newcommand{\N}{\mathbb N}
\DeclareMathOperator{\atom}{atom}
\newcommand{\fset}[2]{\mathcal F^{(#1)}_{#2}}
\title{Tuples as sets}
\author{Adrian Ducourtial}
\date{Source: arXiv:2606.18474v1 (CC BY 4.0)}
\begin{document}
\maketitle

\noindent\emph{Source and licence.} Tuples as sets. Version arXiv:2606.18474v1 (CC BY 4.0), distributed by arXiv under the Creative Commons Attribution 4.0 International licence, http://creativecommons.org/licenses/by/4.0/, as recorded in the arXiv licence metadata for this exact version and shown on the abstract page https://arxiv.org/abs/2606.18474v1. Full licence notice: \texttt{licenses/arxiv-2606.18474-CC-BY-4.0.txt}.\par\smallskip

\noindent\textit{Editorial scope.} Counted unit: the article's theorem thm:diag (source0.tex lines 166--168). The article's complete original proof of it is delivered with it (source0.tex lines 169--222, one \texttt{proof} environment of 54 lines). No mathematics is written, repaired or re-derived: the statement and the proof are the article's own lines, in the article's own notation, and no retained line is substituted or re-flowed.  Retained uncounted as the source-local context the proof needs: lines 61--65; lines 89--91; lines 111--115; lines 137--140; lines 152--154.  Nothing was omitted from any retained span, so the package carries no omission record.  No retained line contains a citation, so the wrapper carries no bibliography and no bibliography entry.  No retained line contains a figure environment, an image-inclusion command or drawing code, so no image is delivered and the package declares no asset dependency.  Editorial formatting only, in addition: an AMS \texttt{amsart} preamble that restates the article's own declarations and macros used by the retained lines, namely the environments \texttt{lemma}, \texttt{theorem} on separate counters, together with the standard packages those lines need.  The retained environments print in this document as Lemma 1, Theorem 1 in order of appearance, whereas the article prints the same items with the numbering of its own full document.  The complete native source of the article as distributed in the arXiv e-print for this exact version is bundled unchanged as \texttt{sources/203081/source0.tex}.
\begin{lemma}\label{lem:var}
    Let $n \geq 0$, $k, \ell \geq 1$, $f \in \fset kn$, and $\pi_{(i)} \in \fset \ell0$ for $1 \leq i \leq k$ be given. The function $f' : D_0^\ell \to D_n$ given by
    \[f'(\mathbf x) = f(\pi_{(1)}(\mathbf x), \dots, \pi_{(k)}(\mathbf x))\]
    is in $\fset \ell n$.
\end{lemma}

\begin{lemma}\label{lem:iota}
    For $n \geq 0$, the function $\iota_n$ is injective.
\end{lemma}

\begin{lemma}\label{lem:permcomm}
    Let $k \geq 1$, $n \geq 0$, and $f \in \fset kn$ be given, and let $\sigma : D_0 \to D_0$ be a permutation. For $\mathbf x = (x_1, \dots, x_k) \in D_0^k$, we have
    \[f(\sigma\mathbf x) = \hat\sigma_n(f(\mathbf x)),\]
    where $\sigma\mathbf x = (\sigma(x_1), \dots, \sigma(x_k))$.
\end{lemma}

\begin{lemma}\label{lem:atom}
    Let $k \geq 1$, $n \geq 0$, and $f \in \fset kn$ be given. For $\mathbf x \in D_0^k$, writing $\mathbf x = (x_1, \dots, x_k)$, we have
    \[\atom_n(f(\mathbf x)) \subseteq \{x_1, \dots, x_k\}.\]
\end{lemma}

\begin{lemma}\label{lem:atomiota}
    We have $\atom_n(\iota_n(x)) = \{x\}$ for all $n \geq 0$.
\end{lemma}

\begin{theorem}\label{thm:diag}
    Let $k, n \in \N$ be given. If $(k, n)$ is good, then $(k+1, n+1)$ is good.
\end{theorem}

\begin{proof}
    Suppose some $f \in \fset kn$ is injective. We have that $f$ is given by $f(\mathbf x) = \{g(\mathbf x) : g \in \mathcal G\}$ for some non-empty $\mathcal G \subseteq \fset k{n-1}$.

    Consider the function $f' : D_0^{k+1} \to D_{n+1}$ given by
    \begin{align*}
        f'(x_1, \dots, x_{k+1})\,=\ &\big\{\left\{g(x_1, \dots, x_k)\right\} : g \in \mathcal G\big\}\ \cup\\
        &\big\{\left\{g(x_1, \dots, x_k), \iota_{n-1}(x_{k+1})\right\} : g \in \mathcal G \big\}.
    \end{align*}
    By use of Lemma~\ref{lem:var}, it can be seen that $f' \in \fset {k+1}{n+1}$. We will show that $f'$ is injective.

    Let $\mathbf x, \mathbf y \in D_0^{k+1}$ be given. Write $\mathbf x = (x_1, \dots, x_{k+1})$, $y = (y_1, \dots, y_{k+1})$, and suppose $f'(\mathbf x) = f'(\mathbf y)$. We wish to show that $\mathbf x = \mathbf y$, i.e., that $x_i = y_i$ for all $1 \leq i \leq k+1$.

    Notice that, for all $g \in \mathcal G$, we have $\{g(x_1, \dots, x_k)\} \in f'(\mathbf x)$, and this set is a singleton. We claim that the converse holds as well: If $s \in f'(\mathbf x)$ is a singleton, then $s = \{g(x_1, \dots, x_k)\}$ for some $g \in \mathcal G$. To see this, note that $s$ is either of the form $\{g(x_1, \dots, x_k)\}$ for some $g \in \mathcal G$ (in which case we are done) or of the form $\{g(x_1, \dots, x_k), \iota_{n-1}(x_{k+1})\}$ for some $g \in \mathcal G$. In the latter case, since $s$ is a singleton, we have $g(x_1, \dots, x_k) = \iota_{n-1}(x_{k+1})$, and hence $s = \{g(x_1, \dots, x_k)\}$ as desired. An analogous argument shows that the singletons in $f'(\mathbf y)$ are exactly the sets in $D_n$ that are of the form $\{g(y_1, \dots, y_k)\}$ for some $g \in \mathcal G$.

    Denoting by $S_{\mathbf x}$ (respectively, $S_{\mathbf y}$) the set of singletons in $f'(\mathbf x)$ (respectively, $f'(\mathbf y)$), we thus have
    \begin{align*}
        \bigcup S_{\mathbf x} &= f(x_1, \dots, x_k),\\
        \bigcup S_{\mathbf y} &= f(y_1, \dots, y_k).
    \end{align*}
    Since $f'(\mathbf x) = f'(\mathbf y)$, we have $S_{\mathbf x} = S_{\mathbf y}$ by definition. Therefore, $\bigcup S_{\mathbf x} = \bigcup S_{\mathbf y}$, i.e., $f(x_1, \dots, x_k) = f(y_1, \dots, y_k)$; and since $f$ is injective, this means that $x_i = y_i$ for $1 \leq i \leq k$. It remains to show that $x_{k+1} = y_{k+1}$.

    Writing $S = S_{\mathbf x} = S_{\mathbf y}$, define
    \begin{align*}
        T_{\mathbf x} &= f'(\mathbf x) \setminus S,\\
        T_{\mathbf y} &= f'(\mathbf y) \setminus S.
    \end{align*}
    Note that the elements of $T_{\mathbf x}$ are exactly those sets in $D_n$ that are of the form $\{g(x_1, \dots, x_k), \iota_{n-1}(x_{k+1})\}$ for some $g \in \mathcal G$ with $g(x_1, \dots, x_k) \neq \iota_{n-1}(x_{k+1})$; similarly, \emph{mutatis mutandis}, for $T_{\mathbf y}$. Moreover, since $f'(\mathbf x) = f'(\mathbf y)$, we have $T_{\mathbf x} = T_{\mathbf y}$. We now split into cases.
    \begin{itemize}
        \item Suppose $T_{\mathbf x} = T_{\mathbf y}$ is empty. This means that for every $g \in \mathcal G$, it is the case that $g(x_1, \dots, x_k) = \iota_{n-1}(x_{k+1})$. Therefore, $f(x_1, \dots, x_k) = \{\iota_{n-1}(x_{k+1})\} = \iota_n(x_{k+1})$. A similar argument shows $f(y_1, \dots, y_k) = \iota_n(y_{k+1})$. Since $f(x_1, \dots, x_k) = f(y_1, \dots, y_k)$, it follows that $\iota_n(x_{k+1}) = \iota_n(y_{k+1})$; hence, by Lemma~\ref{lem:iota}, we have $x_{k+1} = y_{k+1}$ as desired.
        \item Otherwise, let $t \in T_{\mathbf x} = T_{\mathbf y}$ be given. We seek to derive a contradiction. Notice that every member of $T_{\mathbf x}$ contains $\iota_{n-1}(x_{k+1})$; likewise, every member of $T_{\mathbf y}$ contains $\iota_{n-1}(y_{k+1})$. Therefore, in particular, $\iota_{n-1}(x_{k+1}), \iota_{n-1}(y_{k+1}) \in t$. Since $t \notin S$, it must be the case that $t$ has exactly two elements; therefore, $t = \{\iota_{n-1}(x_{k+1}), \iota_{n-1}(y_{k+1})\}$ with $\iota_{n-1}(x_{k+1}) \neq \iota_{n-1}(y_{k+1})$. Notice, however, that $t$ was chosen arbitrarily; and so
        \[T_{\mathbf x} = T_{\mathbf y} = \big\{\left\{\iota_{n-1}(x_{k+1}), \iota_{n-1}(y_{k+1})\right\}\big\}.\]
        This implies, on the one hand, that $\iota_{n-1}(y_{k+1})$ is the unique member of $f(x_1, \dots, x_k)$ that is distinct from $\iota_{n-1}(x_{k+1})$, and, on the other, that $\iota_{n-1}(x_{k+1})$ is the unique member of $f(y_1, \dots, y_k)$ that is distinct from $\iota_{n-1}(y_{k+1})$. As previously established, $f(x_1, \dots, x_k) = f(y_1, \dots, y_k)$, and so together, we have
        \[f(x_1, \dots, x_k) = f(y_1, \dots, y_k) = \{\iota_{n-1}(x_{k+1}), \iota_{n-1}(y_{k+1})\}.\]
        Now, let $\sigma : D_0 \to D_0$ be the permutation $(x_{k+1}\ y_{k+1})$. We have
        \begin{align*}
            f(\sigma(x_1), \dots, \sigma(x_k)) &= \hat\sigma_n(f(x_1, \dots, x_k)) &&\text{by Lemma~\ref{lem:permcomm}}\\
            &= \hat\sigma_n\big(\{\iota_{n-1}(x_{k+1}), \iota_{n-1}(y_{k+1})\}\big)\\
            &= \{\hat\sigma_{n-1}(\iota_{n-1}(x_{k+1})), \hat\sigma_{n-1}(\iota_{n-1}(y_{k+1}))\} &&\text{by def.\ of }\hat\sigma_n\\
            &= \{\iota_{n-1}(\sigma(x_{k+1})), \iota_{n-1}(\sigma(y_{k+1}))\} &&\text{by Lemma~\ref{lem:permcomm}}\\
            &= \{\iota_{n-1}(y_{k+1}), \iota_{n-1}(x_{k+1})\}\\
            &= f(x_1, \dots, x_k).
        \end{align*}
        Since $f$ is injective, this means that $\sigma(x_i) = x_i$ for all $1 \leq i \leq k$.

        Now, because
        \begin{align*}
            \atom_n(f(x_1, \dots, x_k)) &= \atom_n\!\big(\{\iota_{n-1}(x_{k+1}), \iota_{n-1}(y_{k+1})\}\big)\\
            &\supseteq \atom_{n-1}(\iota_{n-1}(x_{k+1}))\\
            &= \{x_{k+1}\} &&\text{by Lemma~\ref{lem:atomiota}},
        \end{align*}
        it follows by Lemma~\ref{lem:atom} that $x_{k+1} \in \{x_1, \dots, x_k\}$, i.e., $x_{k+1} = x_j$ for some $1 \leq j \leq k$; and since $x_j = \sigma(x_j)$, this means $x_{k+1} = \sigma(x_{k+1}) = y_{k+1}$. However, we established earlier that $\iota_{n-1}(x_{k+1}) \neq \iota_{n-1}(y_{k+1})$, implying that $x_{k+1} \neq y_{k+1}$, a contradiction.
    \end{itemize}
    This concludes the proof.
\end{proof}

\end{document}
